\section{Uniform bit complexity}
\label{sec:complexity}

We now charge the constructions to the dense binary input.  The main point
is that neither the degree of the curve nor the order of a primary subgroup
is treated as a constant.  A second issue is the large infinity divisor in
the final label tests: its coefficients are stored in binary, and the
corresponding functions are evaluated from circuits rather than expanded.

Throughout this section, $B=20+(n+1)(L+1)$ is the parameter in
Equation~\eqref{eq:B}.  We use dense polynomial arithmetic and ordinary
Gaussian elimination.  We first count bit operations with indexed arrays,
charging for binary addresses as well as arithmetic, and then allow a
quadratic overhead for a sequential bit implementation.  The exponents
below deliberately allow substantial slack; their purpose is to give one
uniform bound.

\begin{samepage}
\begin{proposition}\label{complexity:bound}
Let $E$ be the largest supplied auxiliary prime, or $2$ when none is needed,
as in Theorem~\ref{thm:reduction}.  The complete algorithm uses
\begin{equation}
 O\left(B^{500000}+B^{100}E^{20}\right)
 \label{complexity:total}
\end{equation}
bit operations.  The same bound applies if the auxiliary primes are found
by upward search and $E$ denotes the largest prime found.  In particular,
if $E=O(B^{20000000})$ with an absolute implied constant, the bit cost is
$O(((n+1)L)^{10^{12}})$.
\end{proposition}
\end{samepage}

\begin{proof}
We first analyze the branch $p>B^{200000}$.  Let $N\le n$ be the degree of
a polynomial passed to the splitting procedure.  In an odd-degree split,
the prime $q\mid N$, the integer $r=v_q(N)$, the extension field $k$, and
the genus $g$ satisfy
\begin{equation}
 \begin{gathered}
 q,N,L\le B,\qquad [k:\mathbf F_p]\le N(q-1)<B^2,\\
 2g<B^2,\qquad m=(q-1)r+1<B^2.
 \end{gathered}
 \label{complexity:basic}
\end{equation}
The intermediate fields used in making one table entry also have degree
at most $B^2$ after the invariant algebra has been extracted.

\smallskip
\noindent\emph{Scalar arithmetic and rational functions.}
A scalar of $k$ has fewer than $B^3$ bits in a basis over $\mathbf F_p$.
A multiplication table, or the equivalent tower representation, gives
addition, multiplication, and inversion by schoolbook arithmetic and
linear algebra.  The algebra $A=k[T]/F$ has dimension $N$ over $k$;
its zero tests with split detection use polynomial gcds, as in
Lemma~\ref{ff:simultaneous}.  Allowing these tests and inversions, we may
charge $O(B^{40})$ bit operations to each scalar instruction over $k$ or
over its simultaneous model $A$.  This bound also covers scalar arithmetic
in the even-degree algebra of dimension $N^2$ over $\mathbf F_p$.
The auxiliary algebras of dimension comparable to $\ell_q$ are charged
separately below.

In calculations over $k(X)$, fractions are reduced after arithmetic
operations.  More specifically, a linear system can first be given one
common polynomial denominator.  For a polynomial matrix of dimension
at most $d$ and entry degrees at most $a$, all its minors have degree at
most $da$.  Pivotal elimination entries are ratios of minors, and solutions
may likewise be obtained by minors after selecting independent rows and
columns.  We use these degree bounds throughout; unreduced expression trees
are never used as representations of rational functions.  Polynomial
determinants can be computed by evaluation and interpolation, with the
grid sizes verified at the end of the proof.

\smallskip
\noindent\emph{One norm equation.}
Each call to $\mathsf{NormSolve}$ from Theorem~\ref{norm:solve} arises
from a divisor of absolute degree
at most $2g$, or from the initial divisor $P_i-\infty_0$.  Compression
expresses its pushforward as $j=h^+/h^-$, with numerator and denominator
degrees at most the absolute degree of that divisor.  Multiplication by
the normalizing constant does not change degrees.  Thus the parameter
$D_G$, the maximum of the numerator and denominator degrees of $G=cj$,
is at most $B^3$.

Here is a size account for the linear construction of global sections of
the order.  The polynomial $H$ is the product of the distinct finite bad
places, and $M=10q^2(N+D_G+1)$ is the local valuation bound used in that
construction.
\begin{center}
\begin{tabular}{ll}
\hline
Quantity & Bound \\
\hline
$D_G$ & $B^3$ \\
$\deg H\le N+2D_G$ & $B^4$ \\
$M$ & $B^7$ \\
Section numerator and common denominator degrees & $B^{12}$ \\
Length of each truncated residue algebra & $B^{13}$ \\
Number of rows and columns in the section system & $B^{17}$ \\
\hline
\end{tabular}
\end{center}
For completeness, the section coordinates have the form
$P_{ab}/H^M$, with $\deg P_{ab}\le M(\deg H+1)$.
There are $q^2$ coordinates.  The local conditions involve $q^2$ matrix
entries at precision at most $3M+3$; summing the residue degrees over all
finite groups gives $\deg H$.  Infinity contributes the same precision
without a residue-degree factor.  These observations give the last three
rows of the table, including a bound on all scalar equations, rather than
only on the number of groups of places.

The root computations in these residue algebras also have bounded sizes.
Their square-free moduli have degree at most $B^4$.  In
Lemma~\ref{ff:unitroot}, every extension degree $d$ is at most $B^4$, so
the bit length of $|k|^d$, and hence the number of primary digits, is at
most $B^8$.  At each digit level there are at most $B^4$ nonempty pieces,
each testing at most $q\le B$ digits.  The kernel-combination search uses
at most $B^9$ scalars.  Thus neither the numerical value of $|k|^d$ nor
the order of its primary subgroup occurs as a loop length.

Let $b$ be the global section lifting the selected rank-one matrix at
infinity.  To form the idempotent $b^q$, retain one common denominator in
the standard cyclic-algebra basis.  Multiplication by $b$ contributes its
denominator and, in reducing powers of the cyclic generator, at most one
additional denominator of $G$.  The numerator degrees grow by at most the
corresponding numerator degrees and a further $O(qN)$ from reducing
powers of $Y$.  Repeating at most $q$ times therefore keeps all degrees
below $B^{15}$.  The system $Ve'=he'$ for $h$ has at most $q$ unknown
coefficients and $q^2$ equations.  The ratio-of-minors bound gives a common
denominator representation
\[
 h=\frac{1}{d_h(X)}\sum_{j=0}^{q-1}a_j(X)Y^j,
 \qquad \deg a_j,\deg d_h\le B^{20}.
\]
The pole estimate for a function in this representation implies
$\|\operatorname{div}h\|\le B^{40}$, where
$\|D\|$ denotes the absolute degree of a divisor.

These sizes also bound the number of operations.  Put $U=B^{25}$.
All polynomial lengths, system dimensions, truncation lengths, and
integer bit lengths just listed are below $U$.  A dense polynomial or
truncated-series operation costs at most $O(U^4)$ scalar operations;
ordinary linear algebra is cubic in its scalar dimensions.  The loops
for residue roots range over groups, degrees, digit levels, nonempty
pieces, possible digits, and powering positions, each with at most $U$
choices.  Constructing local conditions ranges over groups, basis entries,
ansatz coefficients, and precision positions.  Even coefficient-by-
coefficient Hensel lifting and separate construction of every matrix
entry are covered by $O(U^{30})$ scalar and integer operations.
Charging the scalar cost above, and schoolbook costs for the indicated
integers, gives the convenient bound
\begin{equation}
 \text{bit cost of one }\mathsf{NormSolve}\text{ call}=O(B^{2000}).
 \label{complexity:norm}
\end{equation}

\smallskip
\noindent\emph{Divisor arithmetic and reduction.}
After a norm solution in $\mathsf{LambdaDivide}$, the divisor
$E_D=D-\operatorname{div}h$ has absolute degree at most $B^{41}$.
Each orbit partial sum has at most $q$ terms.  The absolute degree of
their coefficientwise maximum is bounded by the sum of their absolute
degrees, so allowing the subtraction of $(\deg V)P_i$ leaves a reduction
input of absolute degree at most $B^{46}$.
All other reduction inputs satisfy the same bound: tracing adds at most
$N$ reduced divisors, and class additions and doublings reduce after
every operation.

Compression gives a bound polynomial of degree at most $\|D\|$ for
each stored divisor.  The Riemann--Roch ansatz for reducing a divisor
of size at most $B^{46}$ consequently uses coefficient and denominator
degrees at most $B^{50}$.  The explicit local pole bounds give series
precision at most $B^{55}$.  Forming the associated principal ideal
also requires an inverse function.  Its multiplication matrix has
dimension $q$ over $k(X)$, so the minor bound gives inverse coefficient
degrees below $B^{54}$.  Clearing denominators for the function and its
inverse, and multiplying the bound polynomials of the ideals being
combined, keeps temporary polynomial degrees below $B^{60}$.
The subspace representation $H^{-1}O/HO$ then has dimension
$2q\deg H<B^{65}$.

The ideal product used for tracing merits a separate check.  Each
component ideal has already been reduced, so its bound polynomial has
degree at most $B^3$ and its generator list has length at most $B^6$.
The determinant formulas multiply at most $N$ such functions.  Before
reduction by $Y^q=F(X)$, the degrees in each of $X$ and $Y$ are less
than $B^{10}$.  The number of scalar parameters in the generator search
is at most $Nb+1$, where $b$ is the generator-list length.  Thus tracing
does not enumerate all choices of one generator from each ideal.

Put $U'=B^{200}$.  The bounds just established place all lengths,
dimensions, precisions, and intermediate generator lists below $U'$,
including the lists of pairwise products in an ideal multiplication.
For an ideal given by generators, closing their span under $X$ and $Y$
requires at most the ambient dimension many strict span enlargements.
Each enlargement uses polynomial arithmetic and linear algebra.
Intersections, inverse-ideal conditions, compression, and the
Riemann--Roch conditions are linear systems in the same representations.
Dense bivariate multiplication takes at most $O((U')^4)$ scalar
operations; determinant interpolation uses at most two grids, each
of length $O((U')^2)$.  Together with the orbit sums and the trace
construction, $O((U')^{30})$ scalar and integer operations bounds any
one of these divisor procedures.

\smallskip
\noindent\emph{Large multiplicities and scalar grids.}
The infinity vectors $W_t$ are the only divisors whose coefficient
magnitudes need not be bounded by a fixed power of $B$.  Their
coefficients have magnitude at most $N(2q)^m$, so their binary lengths
are at most $B^4$, by Equation~\eqref{complexity:basic}.
The construction of $(1-\sigma)D_t-W_t$ uses binary additions and
doublings of divisor classes, reducing the representative after each
instruction and retaining a circuit for the correcting function.
Across all $m$ label tests, at most $B^{10}$ instructions suffice:
there are at most $q$ infinity differences per test and at most $B^4$
binary digits per coefficient.  All leaf functions have the degree
bounds for a single reduction established above.

At a ramification point, evaluate a leaf by its truncated expansion and
propagate its valuation and leading coefficient through the circuit.
Valuations are only added, doubled, or negated.  Even allowing their
magnitudes to double at every circuit instruction, their bit lengths
are below $B^{12}$.  Thus the circuit evaluation does not expand either
a large function or a divisor with a large absolute degree.
Charging $O(B^{100})$ for any of these integer operations is ample.
Together with the bounds on $m$, the number of divisions, and the
primitive procedures above, we obtain
\begin{equation}
 \text{bit cost of one splitting call}=O(B^{20000}).
 \label{complexity:split}
\end{equation}
The even-degree splitter satisfies this bound as well: its algebra has
dimension at most $N^2$, its primary logarithms have at most $L$ digits,
and all exponentiations are binary.

Every scalar grid used above exists in the prime field.  The separator
search in the Berlekamp algebra uses at most $B^3+1$ scalars, the residue
root search at most $B^9$, and the trace search at most $B^8$.
For determinant interpolation, matrix dimensions and entry degrees
are bounded by $U'$, so each grid needs at most $B^{400}+1$ points.
All these lengths are less than $p>B^{200000}$.  The norm solver's
assumption $p>(N+2D_G+1)^2$ follows from the same cutoff and $D_G\le B^3$.
In particular, its derivative-based square-free decompositions have
polynomial degrees smaller than the characteristic.

\smallskip
\noindent\emph{The table and the factorization driver.}
For one auxiliary prime $\ell$, the cyclotomic algebra has dimension
$\ell-1$.  Its automorphisms are computed on the monomial basis, and
their fixed space can be found by stacking at most $\ell$ systems of
at most $\ell$ equations each.  Multiplying a basis of the fixed space
and expressing its products in that basis produces the $q$-dimensional
field representation used subsequently.  Schoolbook polynomial
arithmetic and rectangular elimination, without interpolation grids
depending on $\ell$, give the bound
\begin{equation}
 O(\ell^{10}B^{40})
 \label{complexity:table}
\end{equation}
for this construction, including the other field operations for the
entry but excluding its recursive factorization call.  Trial division
verifies a supplied $\ell$ within the same bound.  If instead it is
found by upward search, trial division and the modular power test for
every integer up to $\ell$ also fit Equation~\eqref{complexity:table}.
Here residues modulo $\ell$ use $O(\log\ell)$ bits and the exponent
test also uses the $L$-bit input $p$.  Powers of $\log\ell$ are absorbed
by the displayed power of $\ell$; no comparison between $\log\ell$
and $\log B$ is needed.  There are at most $B$ table entries, yielding
$O(B^{41}E^{10})$ bit operations outside the recursive calls.

The table is built once, in increasing order of $q$.  Its call to factor
$\Phi_q$ has degree $q-1$ and uses only entries already constructed.
For a factorization input of degree at most $n$, every nontrivial
recursive step decreases degree, and the child degrees in a binary
split sum to the parent degree.  Allowing $O(B^2)$ calls per input,
and including all table inputs, $O(B^4)$ calls is more than sufficient.
The Berlekamp linear systems, polynomial gcds, and evaluations preceding
each split have polynomial dimensions at most those already charged.
Equation~\eqref{complexity:split} therefore bounds all this work by
$O(B^{20004})$.

In the branch $p\le B^{200000}$, no auxiliary table or curve is used.
The extra scalar loop in the Berlekamp splitter has length at most $p$;
all polynomial, matrix, and recursion costs besides this loop are
bounded, for example, by $O(B^{1000})$.  Hence the entire small-prime
branch takes $O(B^{201000})$ bit operations.  The same elementary
operations handle zero derivatives, multiplicities, constants, and
normalization by the leading coefficient.  Multiplicities never
exceed $n$, so their integer bookkeeping is negligible in these bounds.

Combining the two branches gives
\[
 O\left(B^{201000}+B^{41}E^{10}\right)
\]
bit operations with indexed arrays.  The total stored binary data is
bounded by the same quantity; a sequential implementation can locate
entries by scanning the stored data, incurring at most quadratic overhead.
Squaring the preceding bound and enlarging the powers of $B$ gives
Equation~\eqref{complexity:total}.  If $E\le T B^{20000000}$ for a fixed
absolute $T$, its second term is at most $T^{20}B^{400000100}$.
Finally $B=O((n+1)L)$, with an absolute constant, so the stated exponent
$10^{12}$ suffices.
\end{proof}

Proposition~\ref{complexity:bound} completes the uniformity assertion in
Theorem~\ref{thm:reduction}.  It also identifies precisely where a
number-theoretic input can enter the running time: only the numerical
sizes of the auxiliary primes remain outside the polynomial bounds in
$B$ for the algebraic splitting construction.
